\section{A fixed-length theorem in varying dimensions}\label{s:fixed-length}

The same conditional direction law gives a different intermediate
theorem. Instead of testing one continuation on every affine
subsphere in \(\R^d\), we compare learners in different ambient
dimensions. The theorem allows an arbitrary Euclidean output and any
positive error tolerance. A latitude is simulated by a smaller
Gaussian learner, so the induction remains a statement about actual
sample laws.

\begin{proposition}[Dimension-parameterized success bound]\label{s:fixed-proposition}
There are absolute constants \(C\ge2\) and \(d_*\) such that the
following holds for every integer \(d\ge d_*\). Set
\[
 k=p=\lfloor d/32\rfloor,\qquad \gamma=d/32.
\]
Let \(j\ge0\) and \(m\) be integers with \(d/2+2j\le m\le d\).
Consider a fixed-length learner with \(jk\) independent exact Gaussian
samples in dimension \(m\), Borel rules, and at most \(2^{d^2}\)
states at each layer. Its output may be any vector in \(\R^m\).
Its rules may use arbitrary fixed parameters independent of the
target. Its randomized experiment, initialization, any shared rule
seed, and fresh randomness obey the joint measurability, independence,
and information restrictions in Definition~\ref{s:model}, with \(m\)
in place of \(d\). If the target
\(s\) is uniform on \(S^{m-1}\), then for every \(\eta>0\),
\begin{equation}\label{s:fixed-bound}
 \Prob\{\norm{\mathrm{output}-s}\le\eta\}
 \le(C^{j+1}\eta)^\gamma.
\end{equation}
\end{proposition}
\begin{proof}
We induct on \(j\), simultaneously over all stated dimensions,
learners, fixed parameters, and tolerances. All numerical factors introduced below are absolute and independent
of the eventual choice of \(C\). We may condition on the shared seed and start in a
fixed state. The joint independence and measurability requirements
then permit averaging the original conditional probabilities as in
Lemma~\ref{s:borel}.

If \(j=0\), the output is independent of the target. When
\(\eta\ge1\), \eqref{s:fixed-bound} is trivial. When \(\eta<1\),
a zero output cannot succeed. For a nonzero output \(v\), success
implies
\[
 \left\lVert \frac{v}{\norm v}-s\right\rVert
 \le \left|1-\norm v\right|+\norm{v-s}\le2\eta.
\]
Equation~\eqref{s:sharp-cap} therefore bounds the success probability
by \(K_0^d\eta^\gamma\), because \(m-1\ge\gamma\) and
\(\eta<1\). Choosing \(C\) large enough proves the base case,
including random output laws.

Suppose \(j\ge1\). Let \(A\) be the first \(k\) rows, let \(R\) be
their row span, and put
\[
 N_0=R^\perp,\qquad u=P_Rs,\qquad r=\sqrt{1-\norm u^2}.
\]
Almost surely \(u\ne0\) and \(r>0\). Conditional on \(A,u\), the
target is \(u+rz\) with \(z\sim\sigma_{N_0}\), and the first labels
are \(Au\). Thus the first block's selected state supplies no
additional information about \(z\) under this conditioning.

For fixed \(u\), let \(E=u^\perp\) and \(D=m-1\). Write
\(\Prob_{v,s}\) for the law of the continuation from state \(v\)
at target \(s\), and \(\widehat s\) for its output. Define on \(S(E)\)
\[
 f_v(z)=\Prob_{v,u+rz}\{\norm{\widehat s-(u+rz)}\le\eta\}.
\]
These Borel functions use fresh samples and do not depend on the
remaining matrix data. Conditional success is at most
\(\max_v\int f_v\,d\sigma_{N_0}\). By Lemma~\ref{s:residual},
conditional on \(u\), \(N_0\) is a uniform \((m-k)\)-space in
\(E\). Here \(u=\sqrt{1-t}\,e\) fixes the beta fraction and direction
in that lemma. The state maximum is taken before this marginal
subspace law is used.

We obtain a cap bound for \(f=f_v\) from the induction hypothesis.
Fix \(e\in S(E)\), a coordinate \(t\in(-1,1)\), and set
\[
 H=E\cap e^\perp,\qquad
 c=u+rt e,\qquad \rho=r\sqrt{1-t^2}.
\]
The target on this latitude is \(c+\rho z'\) with
\(z'\sim\sigma_H\), where \(\dim H=m-2\) and \(c\perp H\).
Simulate the continuation as a learner for \(z'\) in \(H\).
From a fresh pair \((x',\ip{x'}{z'})\) in \(H\), draw a standard
Gaussian \(x''\) in \(H^\perp\), independently of the joint tuple
consisting of \(z'\), \(x'\), and all previously used learner
randomness. Supply the original
continuation with
\begin{equation}\label{s:latitude-simulation}
 x=x'+x'',\qquad
 y=\ip{x''}{c}+\rho\ip{x'}{z'}=\ip{x}{c+\rho z'}.
\end{equation}
This is exactly a fresh standard Gaussian pair in \(\R^m\) for the
latitude target. The extra Gaussian is used within the transition,
so the state set does not grow. The parameters \(H,c,\rho,v\) are
fixed independently of the varying target \(z'\). At the end, project
the original output onto \(H\) and divide by \(\rho\).
Original error at most \(\eta\) implies simulated error at most
\(\eta/\rho\).

The simulated learner has \((j-1)k\) samples, and
\(m-2\ge d/2+2(j-1)\). The induction hypothesis therefore gives the
latitude average
\[
 \int f(t e+\sqrt{1-t^2}\,z')\,d\sigma_H(z')
 \le\left(\frac{C^j\eta}{r\sqrt{1-t^2}}\right)^\gamma.
\]
Put \(B=C^j\eta/r\). Integrating the latitude bound against the
coordinate density in dimension \(D\), whose remaining exponent is
\((D-3-\gamma)/2\ge0\), gives
\begin{equation}\label{s:fixed-cap}
 \int_{\norm{z-e}\le\tau}f(z)\,d\sigma_E(z)
 \le(c_D/2)B^\gamma\tau^{D-1-\gamma}
 \qquad(\tau>0).
\end{equation}
Indeed, the coordinate interval has length at most \(\tau^2/2\), and
\(1-t^2\le\tau^2\) throughout that interval. At \(\tau=2\), the cap
is the whole sphere. Since \(c_D\le\sqrt d\) and \(D\le d\), both
the cap estimate and this total-mass estimate have coefficients at
most \(4^d\).

We can therefore apply Lemma~\ref{s:cap-moment} to every \(f_v\),
with ambient dimension \(D=m-1\), subspace dimension \(m-k\),
moment \(p=k\), exponent \(a=\gamma\), and
\(\Lambda=B^\gamma\), taking \(A_0=A_1=4\). Its hypotheses hold
for all sufficiently large \(d\):
\(1\le p<m-k\le D\le d\), \(0\le\gamma<D-1\), and
\[
 D-1-\gamma-(p-1)-(k-1)
 =m-\gamma-2k\ge13d/32+2j\ge1.
\]
The maximum conclusion of the lemma, over at most \(2^{d^2}\)
states, yields
\begin{equation}\label{s:fixed-moment}
 \E\left[\max_v\int f_v\,d\sigma_{N_0}\,\middle|\,u\right]
 \le2^{d^2/p}144^d(C^j\eta/r)^\gamma.
\end{equation}
Thus the latitude simulation supplies the cap hypothesis, and the
common moment lemma accounts for all dependence between directions
in the residual space.

It remains to average the inverse radius. By
Lemma~\ref{s:residual},
\(r^2\sim\Beta((m-k)/2,k/2)\). The real-exponent identity
\eqref{s:beta-real} gives
\begin{equation}\label{s:fixed-radius}
 \E r^{-\gamma}
 =\frac{\mathrm B((m-k-\gamma)/2,k/2)}
        {\mathrm B((m-k)/2,k/2)}
 \le4^d.
\end{equation}
To see the bound, note that
\(m-k-\gamma\ge7d/16+2j\), so both exponents in the numerator's
beta integral are nonnegative for large \(d\), and that integral is
at most one. Restricting the denominator to \([1/4,3/4]\) gives
\[
 \mathrm B((m-k)/2,k/2)
 \ge\tfrac12\,4^{-(m/2-2)}\ge4^{-d}.
\]
In particular, this argument applies to the real exponent
\(\gamma=d/32\).

Finally \(p\ge d/64\) for large \(d\). Averaging
\eqref{s:fixed-moment} and applying \eqref{s:fixed-radius} bounds
success by
\[
 (2^{64}\cdot576)^d(C^j\eta)^\gamma.
\]
Since \(\gamma=d/32\), choose the absolute constant \(C\) large
enough that \(C^\gamma\) absorbs this factor and the base-case
factor. This completes the simultaneous induction.
\end{proof}

The dimension condition in the proposition is essential to its
application. Here is the complete reduction that ensures it.
Let \(F=\log(1/\epsilon)\) use natural logarithms, put
\(F_0=8\log C\), and choose an absolute \(a_0>0\) with
\begin{equation}\label{s:fixed-end-constants}
 a_0\le\frac1{8F_0},\qquad a_0\le\frac1{256\log C}.
\end{equation}
Suppose that along arbitrarily large dimensions, learners with
\(M(d)=o(d^2)\) and uniform-prior angular success at least \(2/3\)
satisfy \(T<a_0dF\). Apply Lemma~\ref{s:short-run} with
\(L=F/\log2\), \(c=a_0\log2\), and \(\alpha=2/3\).
It gives \(F=o(d)\), \(T=o(d^2)\), and padded log-width
\(o(d^2)\) on this supposed sequence only.

Pad to \(n=\lceil T/k\rceil=o(d)\) blocks. The padded width
\((T+2)2^M\) is eventually at most \(2^{d^2}\), and
\(d/2+2n\le d\). If \(F\le F_0\), the supposed sample bound and
\eqref{s:fixed-end-constants} give \(T<d/8\), contradicting
\eqref{s:quarter-endpoint}. If \(F>F_0\), then \(k\ge d/64\) and
\[
 (n+1)\log C
 \le(64a_0F+2)\log C\le F/2.
\]
Apply Proposition~\ref{s:fixed-proposition} with \(m=d,j=n\), and
\(\eta=\epsilon\). The original unit-output angular success event
implies Euclidean error at most \(\epsilon\), whereas the proposition
bounds that Euclidean success by
\(\exp(-\gamma F/2)<2/3\) at large \(d\). This contradiction proves
the precision lower bound by the fixed-length route, while retaining
its stronger arbitrary-vector intermediate statement.
